\section{Evaluación del proyecto}
{\consigna
Esta sección aborda el análisis del proyecto con el objetivo de evaluarlo desde dos aspectos: **Técnico** y de **Negocio**.\\
En la **evaluación Técnica** se analizan la factibilidad, los riesgos y todas las cuestiones que aborda el proyecto desde lo tecnológico.\\
En la **evaluación de Negocio** se analizan los aspectos relacionados con la receptividad o posibilidades de adopción del resultado del proyecto por parte de los destinatarios del mismo y las posibilidades de realización desde lo organizacional, legal, económico y demás posibles aspectos no tecnológicos que deban considerarse.
}


\subsection{Tipo de proyecto}
{\consigna
Esta subsección aborda la tipificación del proyecto o su encuadre dentro del tipo de despliegue que el resultado del proyecto va a tener.\\
En clase hemos propuesto los siguientes tipos:
\begin{itemize}
\setlength{\itemsep}{0pt}
  \item Proyectos Masivos
  \item Proyectos Específicos
  \item Proyectos Únicos o a Medida
  \item Proyectos de Innovación o Investigación
\end{itemize}
Esta lista es una propuesta y puede ampliarse o incluso definir una tipificación diferente.
}


El proyecto se clasifica como un Proyecto Masivo por el despliegue previsto de un producto estandarizado, fabricable y distribuible de forma repetible a numerosos usuarios. El resultado mantiene un diseño común compuesto por una aplicación, un conjunto sensor ---que integra el parche sensor descartable, el módulo electrónico reutilizable de lectura y comunicación y la pulsera reutilizable--- y el paquete del sensor. El segmento inicial está constituido por personas con diabetes tipo 2 o prediabetes; la definición de este público no convierte al proyecto en una solución a medida de un único comprador. La configuración individual de parámetros de uso o seguimiento tampoco modifica el diseño común del producto.

Este enfoque exige definir criterios que favorezcan la consistencia entre unidades, una instalación y puesta en marcha comprensibles, la provisión y el recambio de consumibles, el soporte a los usuarios y el análisis de los costos por unidad. Se trata de condiciones a considerar durante el desarrollo: la clasificación no implica que se hayan alcanzado determinados volúmenes o niveles de adopción, ni demuestra por sí misma la factibilidad técnica, la validación del desempeño, la habilitación regulatoria o la madurez comercial del producto.

\subsection{Evaluación de factibilidad técnica}
{\consigna
Esta subsección aborda la evaluación Técnica para analizar la factibilidad, los riesgos y todas las cuestiones que involucran el proyecto desde lo tecnológico.\\
En las diferentes subsecciones se pretenden cubrir todos los aspectos técnicos de la evaluación.
}

Esta sección evalúa la factibilidad técnica del sistema propuesto. Varios de los principios de sensado utilizados —en particular el sensado simultáneo de múltiples biomarcadores en un mismo dispositivo— provienen de desarrollos recientes documentados en literatura científica; el proyecto los adopta y adapta al caso de uso de diabetes, sin que esto implique que la generación de ese conocimiento sea en sí misma el objetivo del trabajo.
\subsubsection{Requerimientos técnicos}
{\consigna
En este apartado deben definirse todos los requerimientos técnicos que deberá cumplir el proyecto, que serán la especificación para el diseño de la solución a construir.\\
Debe contemplar los requisitos que, de cumplirse, cubren las expectativas del proyecto.\\
Todo lo que no se mencione puede no estar y aun así el proyecto ser exitoso.\\
Estas especificaciones suelen usarse durante la verificación para evaluar el cumplimiento del objetivo.\\
Los requerimientos técnicos incluyen los tecnológicos y los funcionales.
}

\noindent\textbf{Requerimientos funcionales:}
\begin{enumerate}
    \item Medir glucosa, lactato, cortisol y pH en sudor de forma no invasiva mediante el parche sensor descartable y el módulo electrónico reutilizable de lectura y comunicación. El pH se utilizará como referencia para la calibración de las demás zonas. Estas mediciones deberán distinguirse de cualquier estimación de glucosa sanguínea.
    \item Medir frecuencia cardíaca, actividad física y temperatura de piel mediante la pulsera, como contexto fisiológico.
    \item Evaluar por separado dos modos propuestos como objetivos de diseño para estimar la evolución de la glucosa: corto plazo (30 a 60 minutos) y horizonte extendido (1 a 8 horas). En ambos modos se deberá indicar el horizonte temporal y la incertidumbre, sin presuponer su viabilidad ni un nivel de precisión ya alcanzado.
    \item Identificar factores probablemente asociados a una variación (alimentación, estrés o actividad física), aclarando que el resultado no establece causalidad clínica.
    \item Presentar información y avisos al usuario acordes con el uso previsto, todavía pendiente de definir, sin emitir recomendaciones automáticas de tratamiento.
    \item Transmitir los datos del módulo electrónico reutilizable y de la pulsera al teléfono por BLE, con una sincronización nominal cada 5 minutos, ajustable según la experiencia de uso. Este intervalo corresponde a la transferencia y no fija la frecuencia de muestreo. Ante una pérdida del enlace, conservar localmente en cada dispositivo al menos 24 horas de datos (estimación: 45 KiB por día y por dispositivo) y transmitir el bloque acumulado al restablecerse BLE.
    \item Definir, antes de evaluar el desempeño, la referencia correspondiente a cada salida, incluidas las mediciones en sudor y cualquier estimación de glucosa sanguínea.
    \item Mantener la solución como complemento de los CGM y documentar como decisión pendiente si los modelos dependerán de datos de un CGM y cómo se realizaría esa integración.
\end{enumerate}
\newpage
\noindent\textbf{Requerimientos técnicos}
\par\medskip
{\small
\renewcommand{\arraystretch}{1.2}
\noindent\begin{tabularx}{\textwidth}{|>{\raggedright\arraybackslash}p{0.28\textwidth}|>{\raggedright\arraybackslash}X|}
\hline
\textbf{Componente} & \textbf{Requerimiento} \\
\hline
Zona glucosa (parche sensor descartable) & Colorimétrica --- reactivo cromogénico específico \\
\hline
Zona lactato (parche sensor descartable) & Colorimétrica --- reactivo cromogénico específico \\
\hline
Zona pH (parche sensor descartable) & Colorimétrica --- referencia/calibración de las demás zonas \\
\hline
Zona cortisol (parche sensor descartable) & Electroquímica --- polímero de impronta molecular (MIP) \\
\hline
Parche sensor descartable & Papel tratado hidrofílica/hidrofóbicamente, plegado tipo origami (microfluídica de papel, bajo costo), incluidas sus zonas con reactivos y su soporte adhesivo. Objetivo de recambio del parche sensor completo cada 7 días de uso, sujeto a validar estabilidad de reactivos, adhesión y desempeño \\
\hline
Lector óptico & Sensor de color para las zonas colorimétricas del parche sensor descartable, integrado en el módulo electrónico reutilizable \\
\hline
Front-end electroquímico & Circuito simple de un solo canal (potenciostato) dedicado a la zona de cortisol, integrado en el módulo electrónico reutilizable \\
\hline
Sensor óptico (pulsera) & PPG para frecuencia cardíaca, idealmente integrado con sensor de temperatura de piel en el mismo IC; pulsera y electrónica reutilizables \\
\hline
IMU (pulsera) & Acelerómetro + giróscopo para detección de actividad/movimiento \\
\hline
Microcontrolador & Bajo consumo, procesa lecturas ópticas y electroquímicas, calibración adaptativa; integrado en el módulo electrónico reutilizable \\
\hline
Comunicación & BLE en el módulo electrónico reutilizable y la pulsera; sincronización nominal cada 5 minutos, ajustable según experiencia de usuario; alcance objetivo de hasta 10 m, incluyendo uso del teléfono en el bolsillo, a validar en pruebas típicas \\
\hline
Almacenamiento local y sincronización & Buffer de al menos 24 horas por dispositivo (estimación: 45 KiB/día). Sincronización cada 5 minutos si BLE está disponible; al reconectar, envío en bloque de los datos acumulados. Para interrupciones superiores a 24 horas no se garantiza conservar el excedente; la aplicación deberá señalar los intervalos faltantes. Validar capacidad y retención ante apagado. \\
\hline
Procesamiento de IA & Modelos a evaluar por separado: estimación temporal (30 a 60 minutos y 1 a 8 horas) e identificación de factores probablemente asociados, sin causalidad clínica. Ejecución local en el teléfono o remota en la nube, coordinada por la aplicación; selección pendiente según costos y requisitos de funcionamiento. La aplicación móvil no requiere reemplazo físico \\
\hline
Alimentación & Batería recargable de bajo consumo en el módulo electrónico reutilizable y en la pulsera reutilizable; autonomía objetivo de al menos 7 días de uso normal para ambos, con perfil de muestreo por definir y a validar \\
\hline
\end{tabularx}
\par}

\subsubsection{Diseño preliminar del producto}
{\consigna
En este apartado debe ir completándose el diseño preliminar. Es un primer bosquejo de diseño que permite tomar dimensión del proyecto y una primera propuesta a alto nivel de solución para cumplir el objetivo y atender los requerimientos.\\
Esta sección lleva más tiempo completarla porque implica la elaboración de un diseño de circuito, sistema, partes, etc.\\
Puede contener diagramas en bloque, especificación de partes y demás detalles que den certidumbre sobre la posibilidad de realización. No es necesario que tenga todo el detalle tecnológico.
}

El sistema se organiza en dos fuentes de datos que convergen en una aplicación central (app/hub), como se muestra en la Figura~\ref{fig:arquitectura-sistema}. El dispositivo del brazo está formado por un parche sensor descartable, ubicado en la parte posterior del brazo, y un módulo electrónico reutilizable de lectura y comunicación que se acopla a él. El parche sensor consiste en papel tratado hidrofílica/hidrofóbicamente y plegado tipo origami, con sus zonas con reactivos y su soporte adhesivo. Sus zonas colorimétricas permiten medir glucosa, lactato y pH en sudor mediante reactivos cromogénicos y un sensor de color situado en el módulo electrónico reutilizable; la zona de pH se utiliza como referencia para la calibración de las demás zonas. Una zona electroquímica separada del parche sensor, basada en un polímero de impronta molecular (MIP), permite medir cortisol en sudor mediante un potenciostato de un solo canal integrado en el módulo electrónico reutilizable. Estas mediciones no constituyen por sí mismas mediciones de glucosa sanguínea: cualquier salida que la estime deberá especificar su referencia de evaluación. La pulsera reutilizable, en la muñeca, integra el sensor óptico PPG, el sensor de temperatura, la IMU y su propio módulo BLE, todos reutilizables, aportando el contexto fisiológico de frecuencia cardíaca, actividad física y temperatura de piel. El módulo electrónico reutilizable y la pulsera transmiten sus datos por BLE a la aplicación central (hub), donde se proponen modelos de inteligencia artificial que deberán evaluarse por separado: un modelo temporal con objetivos de diseño de corto plazo (30 a 60 minutos) y horizonte extendido (1 a 8 horas), sin presuponer la viabilidad ni la precisión de ninguno de ellos; y un modelo que identifica factores probablemente asociados a la variación al relacionar biomarcadores y señales fisiológicas de contexto, sin establecer causalidad clínica. Cada salida deberá presentar su horizonte temporal, incertidumbre y referencia de evaluación según corresponda. La aplicación presentará información y avisos coherentes con el uso previsto, aún pendiente de cerrar, sin recomendaciones automáticas de tratamiento. El sistema se concibe como complemento de los CGM; queda pendiente determinar si los modelos dependerán de sus datos y cómo se integrarían.

La aplicación central se ejecutará en el teléfono del usuario y recibirá los datos del módulo electrónico reutilizable y la pulsera. Para ejecutar los modelos de IA se evaluarán dos alternativas: procesamiento local en el teléfono o procesamiento remoto en un servicio en la nube, al que la aplicación enviaría los datos necesarios y del que recibiría los resultados. La elección se realizará durante el diseño detallado, comparando los costos de implementación y operación, la capacidad de procesamiento del teléfono, el tiempo de respuesta, la disponibilidad de conexión a Internet y los requisitos de privacidad y protección de datos. No se compromete en esta etapa la implementación de ambas modalidades. Si se elige la nube, deberá definirse el comportamiento del sistema ante la pérdida de conexión, indicando al usuario qué funciones no estarán disponibles.

La comunicación entre el módulo electrónico reutilizable, la pulsera y el teléfono tendrá un intervalo nominal de sincronización de 5 minutos, ajustable según la experiencia de uso. Se buscará un alcance BLE que cubra el uso del teléfono en el bolsillo y hasta 10 m en condiciones típicas; ambos aspectos deberán validarse en pruebas. Este intervalo es de transferencia y no define la frecuencia de muestreo. Si el teléfono se aleja o se apaga mientras el módulo y la pulsera siguen encendidos, estos conservarán localmente hasta 24 horas de datos cada uno (estimación: 45 KiB por día y dispositivo) y enviarán el bloque acumulado al restablecerse BLE. Si la interrupción supera las 24 horas, no se garantiza conservar el excedente y la aplicación deberá indicar los intervalos faltantes, sin presentar la serie como continua. Si los propios dispositivos se apagan, no se adquieren mediciones durante ese período; la persistencia de los datos ya almacenados ante un apagado dependerá de la memoria elegida y deberá verificarse. La estimación de memoria se considera compatible con microcontroladores habituales, sujeta a validación con el formato, las marcas temporales, los metadatos y las frecuencias de muestreo definidos.

El consumible descartable del sistema será el parche sensor de papel completo, incluidas sus zonas con reactivos y su soporte adhesivo de fijación a la piel del brazo. Se establece como objetivo reemplazar el parche sensor completo cada 7 días de uso, sujeto a validar la estabilidad de los reactivos, la adhesión y el desempeño, conservando el módulo electrónico reutilizable de lectura y comunicación que se acopla a él. La pulsera y su electrónica también serán reutilizables. La aplicación móvil no requiere reemplazo físico. La vida útil objetivo de 7 días y las instrucciones para retirar, reemplazar y volver a acoplar el parche sensor deberán validarse durante el desarrollo.

En la alternativa de procesamiento remoto, los datos llegarían a la nube mediante el acceso a Internet del teléfono (Wi-Fi o datos móviles). La aplicación enviaría los datos necesarios y recibiría los resultados por esa misma conexión. El módulo electrónico reutilizable y la pulsera se comunicarían con el teléfono por BLE, sin conectarse directamente a la nube.

\begin{figure}[H]
    \centering
    \resizebox{\linewidth}{!}{%
    \begin{tikzpicture}
    \tikzset{
        bloque/.style={draw=tpazul!65,fill=tpclaro,rounded corners=3pt,
            text width=2.5cm,minimum height=1.15cm,inner sep=5pt,
            align=center,font=\sffamily\small,execute at begin node={\hyphenpenalty=10000}},
        sensor/.style={bloque,fill=tpverde},
        energia/.style={bloque,fill=tpenergia},
        externo/.style={bloque,fill=white,draw=tpazul!40},
        flujo/.style={-{Stealth[length=2mm]},line width=.8pt,draw=tpazul},
        potencia/.style={flujo,dashed,draw=orange!70!black},
        relacion/.style={line width=.8pt,draw=tpazul},
        rotulo/.style={font=\sffamily\scriptsize,fill=white,inner sep=2pt,text=tpazul},
        marco/.style={draw=tpazul!35,rounded corners=5pt,fill=black!1,inner sep=12pt},
        titulo/.style={font=\sffamily\bfseries\large,text=tpazul,anchor=west}
    }
    \node[titulo] at (-1.5,4.4) {Parche descartable};
    \node[externo] (sudor) at (0,0) {Sudor\\de la piel};
    \node[sensor] (cap) at (3.6,0) {Captación\\y distribución\\en papel};
    \node[sensor] (glu) at (8,2.7) {Zona de glucosa\\Reactivo};
    \node[sensor] (lac) at (8,.9) {Zona de lactato\\Reactivo};
    \node[sensor] (ph) at (8,-.9) {Zona de pH\\Indicador};
    \node[sensor] (cor) at (8,-2.7) {Zona de cortisol\\Electrodos y MIP};
    \node[externo] (opt) at (13.1,.9) {Lectura óptica\\por el lector};
    \node[externo] (ele) at (13.1,-2.7) {Lectura eléctrica\\por el lector};
    \draw[flujo] (sudor)--(cap);
    \draw[relacion] (cap.east)--(5.6,0);
    \draw[relacion] (5.6,-2.7)--(5.6,2.7);
    \foreach \n in {glu,lac,ph,cor} {\draw[flujo] (5.6,0|-\n.west)--(\n.west);}
    \draw[relacion] (glu.east)--(10.1,2.7)--(10.1,-.9)--(ph.east);
    \draw[relacion] (lac.east)--(10.1,.9);
    \draw[flujo] (10.1,.9)--node[rotulo,above]{Color}(opt.west);
    \draw[flujo] (cor.east)--node[rotulo,above]{Señal eléctrica}(ele.west);
    \node[font=\sffamily\small,text=tpazul] (base) at (5.8,-4.1) {Soporte y adhesivo en contacto con la piel};
    \begin{scope}[on background layer]
        \node[marco,fit=(cap)(glu)(lac)(ph)(cor)(base)] {};
    \end{scope}
    \end{tikzpicture}}
    \caption{Esquema funcional del parche sensor descartable: captación y distribución del sudor, zonas de sensado y lecturas óptica y electroquímica realizadas por el módulo reutilizable.}
    \label{fig:parche-descartable}
\end{figure}



\begin{center}
\resizebox{\linewidth}{!}{%
\begin{tikzpicture}
\node[titulo] at (-1.5,2.4) {Lector reutilizable del brazo};
\node[externo] (color) at (0,.9) {Zonas de color\\del parche};
\node[externo] (cort) at (0,-.9) {Electrodos\\del parche};
\node[bloque] (optica) at (4,.9) {Iluminación\\y sensor de color};
\node[bloque] (pot) at (4,-.9) {Potenciostato\\Lectura eléctrica};
\node[bloque] (mc) at (8.1,0) {Microcontrolador\\Adquisición\\y calibración};
\node[bloque,text width=1.8cm] (ble) at (11.8,0) {Enlace\\BLE};
\node[externo,text width=1.6cm] (app) at (15,0) {Aplicación\\móvil};
\node[energia,text width=6cm] (pwr) at (7.7,-3.2) {Batería recargable\\Carga, protección y regulación};
\draw[flujo] (color)--(optica);
\draw[flujo,<->] (cort)--(pot);
\draw[flujo] (optica.east)--++(.5,0)|-([yshift=.27cm]mc.west);
\draw[flujo] (pot.east)--++(.5,0)|-([yshift=-.27cm]mc.west);
\draw[flujo] (mc)--(ble);
\draw[flujo] (ble)--(app);
\draw[potencia] (pwr.north)--(pwr.north|-mc.south);
\draw[potencia] (pwr.west)--(2.2,-3.2)--(2.2,.45)--([yshift=-.45cm]optica.west);
\draw[potencia] (2.2,-1.35)--([yshift=-.45cm]pot.west);
\draw[potencia] (pwr.east)-|(ble.south);
\node[rotulo] at (10.5,-2.05) {Energía};
\node[font=\sffamily\small,text=tpazul] (base) at (7.5,-4.4) {Carcasa opaca y acople al parche};
\begin{scope}[on background layer]
\node[marco,fit=(optica)(pot)(mc)(ble)(pwr)(base)] {};
\end{scope}
\end{tikzpicture}}
\end{center}


\begin{center}
\resizebox{\linewidth}{!}{%
\begin{tikzpicture}
\node[titulo] at (-1.5,4.3) {Pulsera reutilizable};
\node[sensor] (ppg) at (0,2.5) {Sensor óptico PPG\\Frecuencia cardíaca};
\node[sensor] (imu) at (0,.5) {Acelerómetro\\y giróscopo};
\node[sensor] (tmp) at (0,-1.5) {Temperatura\\de piel};
\node[bloque] (mc) at (5.2,.5) {Microcontrolador\\Adquisición\\y filtrado};
\node[bloque] (ble) at (9.1,.5) {Comunicación\\BLE};
\node[externo] (app) at (13,.5) {Aplicación\\móvil};
\node[energia,text width=5.5cm] (pwr) at (5.5,-3.6) {Batería recargable\\Carga, protección y regulación};
\draw[flujo] (ppg.east)--(2.3,2.5)|-([yshift=.38cm]mc.west);
\draw[flujo] (imu)--(mc);
\draw[flujo] (tmp.east)--(2.3,-1.5)|-([yshift=-.38cm]mc.west);
\draw[flujo] (mc)--(ble);
\draw[flujo] (ble)--(app);
\draw[potencia] (pwr.north)--(pwr.north|-mc.south);
\draw[potencia] (pwr.east)-|(ble.south);
\draw[potencia] (pwr.west)--(-2.1,-3.6)--(-2.1,2.5)--(ppg.west);
\draw[potencia] (-2.1,.5)--(imu.west);
\draw[potencia] (-2.1,-1.5)--(tmp.west);
\node[rotulo] at (6.7,-2.15) {Energía};
\node[font=\sffamily\small,text=tpazul] (base) at (4.3,-4.85) {Carcasa, contacto con la piel y correa};
\begin{scope}[on background layer]
\node[marco,fit=(ppg)(imu)(tmp)(mc)(ble)(pwr)(base),inner xsep=22pt] {};
\end{scope}
\end{tikzpicture}}
\end{center}



\begin{center}
\resizebox{\linewidth}{!}{%
\begin{tikzpicture}
\node[titulo] at (-1.5,2.2) {Aplicación móvil};
\node[externo] (in) at (0,0) {Datos del lector\\y de la pulsera};
\node[bloque] (rx) at (3.9,0) {Recepción\\por BLE};
\node[bloque] (prep) at (7.8,0) {Sincronización\\y control de calidad};
\node[bloque] (model) at (11.7,0) {Análisis\\y estimaciones};
\node[bloque] (hist) at (7.8,-2.5) {Historial\\de mediciones};
\node[bloque] (ui) at (11.7,-2.5) {Visualización\\y avisos};
\node[externo] (usr) at (15.6,-2.5) {Usuario};
\draw[flujo] (in)--(rx);
\draw[flujo] (rx)--(prep);
\draw[flujo] (prep)--(model);
\draw[flujo] (prep)--(hist);
\draw[flujo] (hist.east)--++(.5,0)|-([yshift=-.4cm]model.west);
\draw[flujo] (model)--(ui);
\draw[flujo] (hist.south)--++(0,-.65)-|(ui.south);
\draw[flujo] (ui)--(usr);
\begin{scope}[on background layer]
\node[marco,fit=(rx)(prep)(model)(hist)(ui),inner ysep=24pt] {};
\end{scope}
\end{tikzpicture}}
\end{center}

\begin{center}
\resizebox{\linewidth}{!}{%
\begin{tikzpicture}
\node[titulo] at (-1.6,2.3) {Caja del producto};
\node[bloque,text width=3.2cm] (caja) at (5.5,.3) {Caja exterior};
\node[bloque,text width=3.2cm] (guia) at (0,-2.3) {Guía de inicio\\y manual};
\node[bloque,text width=3.2cm] (proteccion) at (5.5,-2.3) {Protección interna\\y alojamientos};
\node[bloque,text width=3.2cm] (etiqueta) at (11,-2.3) {Etiquetas\\e identificación};
\node[externo,text width=3.2cm] (equipos) at (0,-5.4) {Lector y pulsera};
\node[externo,text width=3.2cm] (parche) at (5.5,-5.4) {Parche\\en sobre individual};
\node[externo,text width=3.2cm] (accesorios) at (11,-5.4) {Accesorios\\de carga};
\draw[relacion] (caja.south)--(5.5,-1.05);
\draw[relacion] (0,-1.05)--(11,-1.05);
\draw[relacion] (0,-1.05)--(guia.north);
\draw[relacion] (5.5,-1.05)--node[rotulo,right]{Contiene}(proteccion.north);
\draw[relacion] (11,-1.05)--(etiqueta.north);
\draw[relacion] (proteccion.south)--(5.5,-3.8);
\draw[relacion] (0,-3.8)--(11,-3.8);
\draw[relacion] (0,-3.8)--(equipos.north);
\draw[relacion] (5.5,-3.8)--node[rotulo,right]{Aloja y protege}(parche.north);
\draw[relacion] (11,-3.8)--(accesorios.north);
\end{tikzpicture}}
\end{center}



\begin{figure}[htbp]
    \centering
    \begin{tikzpicture}[
        bloque/.style={draw=black!65, rounded corners=3pt, fill=gray!5,
            align=center, font=\small, text width=4.3cm,
            minimum height=2.4cm, inner sep=8pt},
        flujo/.style={-{Stealth[length=3mm]}, thick, draw=black!70},
        remoto/.style={{Stealth[length=3mm]}-{Stealth[length=3mm]},
            thick, dashed, draw=black!70},
        etiqueta/.style={font=\footnotesize, align=center, fill=white, inner sep=3pt}
    ]
        \node[bloque] (parche) at (-3.5,0) {\textbf{Dispositivo del brazo}\\Parche sensor descartable\\(papel, reactivos y adhesivo)\\Glucosa, lactato, cortisol; pH\\Módulo electrónico reutilizable\\Lectura y comunicación BLE};
        \node[bloque] (pulsera) at (3.5,0) {\textbf{Pulsera (muñeca)}\\Pulsera y electrónica\\reutilizables\\Frecuencia cardíaca,\\actividad y temperatura};
        \node[bloque] (hub) at (-3.5,-4.8) {\textbf{App en el teléfono}\\Recepción de datos\\IA local si se elige\\procesamiento en el celular};
        \node[bloque, dashed] (nube) at (3.5,-4.8) {\textbf{Servicio en la nube}\\IA remota si se elige\\procesamiento en la nube\\Alternativa por definir};
        \node[bloque, minimum height=1.6cm] (alertas) at (-3.5,-8) {\textbf{Información y avisos}\\Presentados en el teléfono\\Según uso previsto por definir};
        \draw[flujo] (parche.south) -- node[etiqueta, right] {BLE} (-3.5,-2)
            -| ([xshift=-1.2cm]hub.north);
        \draw[flujo] (pulsera.south) -- node[etiqueta, right] {BLE} (3.5,-2)
            -| ([xshift=1.2cm]hub.north);
        \draw[remoto] (hub.east) -- node[etiqueta, above] {Internet}
            node[etiqueta, below] {Datos /\\resultados} (nube.west);
        \draw[flujo] (hub.south) -- (alertas.north);
    \end{tikzpicture}
    \caption{Flujo de datos del sistema. El dispositivo del brazo comprende el parche sensor descartable completo (papel, zonas con reactivos y soporte adhesivo) y el módulo electrónico reutilizable de lectura y comunicación que se acopla a él. El módulo electrónico y la pulsera, cuya electrónica también es reutilizable, envían sus datos al teléfono por BLE. La aplicación móvil no requiere reemplazo físico; el recambio semanal del parche sensor es un objetivo pendiente de validar. El bloque y la conexión discontinuos representan la alternativa de IA en la nube, accesible mediante Internet (Wi-Fi o datos móviles); si se elige IA local, el procesamiento se realiza en el teléfono sin ese recorrido. Los modelos estimarán la evolución de la glucosa (30--60 min y 1--8 h) e identificarán factores probablemente asociados. Las mediciones del parche sensor descartable corresponden a sudor; la referencia de evaluación de cada salida deberá definirse.}
    \label{fig:arquitectura-sistema}
\end{figure}

    \paragraph{Análisis AMFE (preliminar)}
    \label{sec:amfe-preliminar}
    {\consigna
    Aquí debe incorporarse el Análisis AMFE preliminar e ir enriqueciéndolo a lo largo del diseño y desarrollo del proyecto.
    }


{\footnotesize
\renewcommand{\arraystretch}{1.2}
\begin{longtable}{|>{\raggedright\arraybackslash}p{\dimexpr0.18\textwidth-2\tabcolsep-1.2\arrayrulewidth\relax}|>{\raggedright\arraybackslash}p{\dimexpr0.22\textwidth-2\tabcolsep-1.2\arrayrulewidth\relax}|>{\raggedright\arraybackslash}p{\dimexpr0.18\textwidth-2\tabcolsep-1.2\arrayrulewidth\relax}|>{\raggedright\arraybackslash}p{\dimexpr0.13\textwidth-2\tabcolsep-1.2\arrayrulewidth\relax}|>{\raggedright\arraybackslash}p{\dimexpr0.29\textwidth-2\tabcolsep-1.2\arrayrulewidth\relax}|}
\hline
\textbf{Componente / función} & \textbf{Modo de falla} & \textbf{Efecto} & \textbf{Severidad} & \textbf{Mitigación propuesta} \\
\hline
\endfirsthead
\hline
\textbf{Componente / función} & \textbf{Modo de falla} & \textbf{Efecto} & \textbf{Severidad} & \textbf{Mitigación propuesta} \\
\hline
\endhead
Zonas colorimétricas (glucosa, lactato, pH) & Interferencia por condiciones de luz ambiente al leer el color & Medición imprecisa & Media & Lector con carcasa que aísla luz ambiente, o calibración por referencia de color interna \\
\hline
Zonas colorimétricas & Degradación del reactivo cromogénico por humedad o tiempo de almacenamiento & Lectura nula o errónea & Media & Empaque sellado, indicación de vida útil, control de lote en fabricación \\
\hline
Secuenciado de reacciones (plegado de papel) & Plegado incorrecto por parte del usuario al colocar el parche sensor descartable & Mezcla de zonas, lectura cruzada & Media & Diseño de plegado guiado mecánicamente (una sola forma posible de doblar) \\
\hline
Zona cortisol (MIP electroquímico) & Baja especificidad, interferencia de otras moléculas del sudor & Asociación incorrecta con estrés & Media & Evaluar especificidad y contrastar con señales de contexto; no presentar el resultado como diagnóstico ni como causalidad clínica \\
\hline
Comunicación BLE & Pérdida de conexión entre dispositivo y hub & Interrupción de datos o aviso no emitido & Media & Sincronización cada 5 min y alcance objetivo de hasta 10 m (incluido teléfono en bolsillo), a validar; buffer local de 24 h (45 KiB/día por dispositivo) y envío en bloque al reconectar; advertir intervalos faltantes si se excede el buffer \\
\hline
Batería (módulo electrónico reutilizable o pulsera) & Descarga antes de lo esperado & Vacío de datos, pérdida de contexto & Media & Diseño de bajo consumo e indicador de batería baja; verificar autonomía objetivo de 7 días en el perfil nominal de uso \\
\hline
Modelo de factores asociados (IA) & Asociación incorrecta entre una variación y un factor de contexto & Información o aviso inadecuado al usuario & Alta & Presentar factores como probablemente asociados, informar incertidumbre y evitar inferencias causales o recomendaciones automáticas de tratamiento \\
\hline
\end{longtable}
}

\subsubsection{Inventario de los diferentes aspectos a cubrir}
{\consigna
En este apartado debe enumerarse y describirse todos los aspectos involucrados en el proyecto que son necesarios para cumplir el objetivo.\\
Dependiendo del objetivo, será la cantidad de aspectos a tener en cuenta.\\
Se mencionan algunos ejemplos que podrían aplicar:

\begin{itemize}
\setlength{\itemsep}{3 pt}
    \item Línea de Producción / fabricación
    \item QA
    \item Soporte
    \item Repuestos, Entrenamiento
    \item Servicios Especializados
\end{itemize}
}


\begin{itemize}
    \item \textbf{Línea de producción/fabricación:} fabricación de parches sensores descartables de papel plegado (microfluídica de papel), incluidas sus zonas con reactivos y su soporte adhesivo, y ensamblado del módulo electrónico reutilizable de lectura óptica/electroquímica y comunicación y de la pulsera reutilizable.
    \item \textbf{QA:} control de calidad por lote de reactivos cromogénicos, calibración de color de cada lector, calibración de cada unidad de pulsera.
    \item \textbf{Soporte:} atención a usuarios ante fallas de sincronización, dudas sobre plegado y recambio del parche sensor descartable, acoplamiento del módulo electrónico reutilizable o interpretación de la información y los avisos; mantenimiento del módulo electrónico y de la pulsera reutilizables y actualización de firmware remota.
    \item \textbf{Repuestos y entrenamiento:} provisión y recambio del parche sensor descartable completo (papel, zonas con reactivos y soporte adhesivo), con recambio objetivo cada 7 días, sujeto a validar la estabilidad de los reactivos, la adhesión y el desempeño; conservación y mantenimiento del módulo electrónico reutilizable y de la pulsera con su electrónica reutilizable. La aplicación móvil no requiere reemplazo físico. Capacitación a usuarios y profesionales para efectuar el recambio del parche sensor y el acoplamiento del módulo electrónico, e interpretar las mediciones, estimaciones, incertidumbres y avisos dentro del uso previsto.
    \item \textbf{Servicios especializados:} laboratorio de química/bioquímica para seleccionar reactivos cromogénicos existentes, adaptar su formulación e inmovilización en el soporte de papel y validar su desempeño en sudor; para la zona de cortisol, preparación e integración del MIP a partir de protocolos publicados, con evaluación de selectividad y estabilidad. No se plantea desarrollar nuevas químicas de detección desde cero, sino adaptar soluciones existentes al dispositivo propuesto. 
\end{itemize}

\subsection{Evaluación de factibilidad de negocio}
{\consigna
Esta sección aborda la evaluación de Negocio, analizando los aspectos relacionados con la receptividad o posibilidades de adopción del resultado del proyecto por parte de los destinatarios del mismo y las posibilidades de realización desde lo organizacional, legal, económico y demás cuestiones no tecnológicas que deban considerarse.
}

\subsubsection{Análisis FODA}
{\consigna
Esta sección incorpora un análisis FODA del proyecto (del resultado del proyecto).\\
Esta sección, como todas, es dinámica y se enriquece a lo largo del desarrollo del proyecto.
}

El análisis considera el alcance conceptual actual del proyecto. Las fortalezas se plantean como ventajas internas esperadas del producto. Las oportunidades corresponden a condiciones externas que podrían favorecer su desarrollo o adopción.

\noindent\textbf{Fortalezas}
\begin{itemize}
    \item Mayor apoyo para interpretar las variaciones de glucosa frente a la consulta de un valor o una predicción aislada: el producto ofrece información sobre factores probablemente asociados, como alimentación, estrés o actividad física, para que el usuario no tenga que interpretar únicamente un número por su cuenta. Esta información complementa el seguimiento con CGM, sin constituir un diagnóstico de la causa.
    \item Menor esfuerzo de registro frente a la transcripción manual de mediciones: la recepción automática de datos del módulo electrónico reutilizable y la pulsera y su organización en un historial común evitan que el usuario copie y reúna esas mediciones por su cuenta para consultar su evolución.
    \item Seguimiento de varios biomarcadores en sudor sin punciones ni extracción de sangre para obtener esas mediciones, como alternativa no invasiva para reunir información metabólica adicional. Esta ventaja no implica sustituir los controles sanguíneos ni el CGM.
    \item Menor necesidad de recomprar componentes al reemplazar el parche sensor descartable: el usuario puede conservar el módulo electrónico reutilizable de lectura y comunicación y la pulsera reutilizable, en lugar de desechar el conjunto completo cada vez que reemplaza el parche sensor de papel con sus reactivos y adhesivo.
\end{itemize}

\noindent\textbf{Oportunidades}
\begin{itemize}
    \item Posibilidad de atender a personas con diabetes tipo 2 o prediabetes que busquen comprender mejor la relación entre sus hábitos cotidianos y las variaciones de glucosa, mediante una propuesta de seguimiento complementario y personalizado.
    \item Posibilidad de establecer alianzas con clínicas y profesionales de salud interesados en evaluar herramientas de seguimiento entre consultas, validar la utilidad de la información contextual y definir cómo presentar los avisos.
    \item Interés en el uso de biomarcadores para el monitoreo de enfermedades crónicas.
\end{itemize}

\noindent\textbf{Debilidades}
\begin{itemize}
    \item Experiencia y recursos especializados limitados: el equipo no cuenta actualmente con experiencia práctica demostrada, antecedentes propios ni infraestructura especializada para desarrollar y validar sensores bioquímicos, dispositivos médicos y modelos de inteligencia artificial aplicados a información clínica. También carece de experiencia directa en validación clínica y procesos regulatorios, por lo que dependerá de laboratorios, profesionales de salud e instituciones externas para cubrir estas capacidades.
    \item Experiencia limitada del equipo en el desarrollo integral de dispositivos médicos: como estudiantes de Ingeniería Electrónica, los integrantes cuentan con formación teórica en electrónica, pero aún deben adquirir experiencia práctica en su aplicación al diseño de dispositivos médicos. También será necesario ampliar conocimientos en biosensores, modelos de IA, validación clínica, regulación e industrialización, mediante aprendizaje específico, asesoramiento especializado y alianzas con profesionales de esas áreas.
    \item Desempeño aún no demostrado: el proyecto no cuenta con un prototipo integrado ni validación propia que permita comprobar la calidad de las mediciones y de las predicciones propuestas.
    \item La colocación del parche sensor descartable, el acoplamiento del módulo electrónico reutilizable y el uso conjunto con la pulsera pueden resultar complejos para el usuario, lo que podría dificultar la utilización adecuada del sistema.
\end{itemize}

\noindent\textbf{Amenazas}
\begin{itemize}
    \item Disponibilidad o acceso insuficiente a datos adecuados para la IA: la falta de registros sincronizados de biomarcadores, señales fisiológicas y eventos de referencia podría limitar el entrenamiento y la evaluación de los modelos, comprometiendo la confiabilidad de la información ofrecida por el producto. Si se determina que los modelos requieren datos de un CGM, las restricciones de acceso a esos registros podrían agravar esta limitación.
    \item Riesgo de dificultades en la integración y calibración: la coordinación de lecturas colorimétricas, electroquímicas y fisiológicas podría no alcanzar la consistencia requerida, comprometiendo el desempeño del sistema y aumentando los plazos y costos de desarrollo.
    \item Fabricantes de CGM u otros desarrolladores podrían incorporar funciones de interpretación contextual o asociación de eventos antes de que el proyecto alcance una versión comercial, reduciendo su diferenciación.
    \item Las exigencias de evidencia, seguridad y autorización que resulten aplicables al uso previsto podrían extender los plazos y aumentar los costos del desarrollo. La clasificación regulatoria deberá confirmarse y no se considera resuelta por utilizar una tecnología de sensado determinada.
    \item Usuarios y posibles financiadores de la compra, como obras sociales o aseguradoras, podrían no considerar suficiente el beneficio adicional frente al costo del sistema y al recambio de parches sensores descartables, con recambio objetivo cada 7 días, sujeto a validar su estabilidad, adhesión y desempeño, limitando su adopción incluso si la solución resulta técnicamente viable.
    \item La falta de acceso a financiamiento externo o a acuerdos con instituciones para realizar ensayos y obtener datos podría demorar el paso del diseño conceptual a un prototipo validado.
\end{itemize}

\subsubsection{Análisis de Riesgos}
{\consigna
En esta sección se deben identificar los riesgos del proyecto, es decir, todos los factores que puedan atentar contra el éxito del mismo.\\
Hemos dividido el análisis en dos partes: Matriz de Riesgos y Análisis AMFE.
}

Este análisis identifica los riesgos que pueden afectar la elaboración de los requerimientos, el diseño preliminar de hardware, la arquitectura de IA y la evaluación de factibilidad técnica y de negocio del proyecto. Se consideran sus efectos sobre los plazos, el alcance y la calidad de los entregables, mientras que los modos de falla del producto se abordan en el AMFE de la sección~\ref{sec:amfe-preliminar}.

\noindent\textbf{a) Identificación de los riesgos}
\begin{itemize}
    \item R1 $\rightarrow$ Evidencia insuficiente para sustentar la factibilidad del sensado propuesto.
    \item R2 $\rightarrow$ Falta de información para especificar el entrenamiento y la evaluación de la IA.
    \item R3 $\rightarrow$ Subestimación del esfuerzo y demoras en la elaboración de los entregables.
    \item R4 $\rightarrow$ Cambios tardíos o ampliación del alcance durante el desarrollo del trabajo.
    \item R5 $\rightarrow$ Inconsistencias entre el diseño técnico y las conclusiones de factibilidad y negocio.
\end{itemize}

\noindent\textbf{b) Estimación de las consecuencias de los riesgos}

Para este análisis se propone una escala ordinal de 1 a 10: la severidad ($S$) expresa el impacto sobre la calidad, el alcance o los plazos de los entregables, y la ocurrencia ($O$) expresa la posibilidad estimada de que el riesgo se materialice. Un valor mayor indica mayor impacto o posibilidad de ocurrencia; $O$ no representa un porcentaje. Las puntuaciones son estimaciones preliminares que deberán revisarse con el equipo, no frecuencias medidas.

\noindent R1 $\rightarrow$ Evidencia insuficiente para sustentar la factibilidad del sensado propuesto.
\begin{itemize}
    \item Severidad ($S$): 9, porque podría obligar a revisar los requerimientos y la arquitectura del parche o impedir una conclusión fundamentada sobre su factibilidad.
    \item Ocurrencia ($O$): 6, porque todavía deben contrastarse los rangos, interferencias, estabilidad de los reactivos y condiciones de uso de las zonas colorimétricas y del MIP de cortisol.
\end{itemize}

\noindent R2 $\rightarrow$ Falta de información para especificar el entrenamiento y la evaluación de la IA.
\begin{itemize}
    \item Severidad ($S$): 8, porque dejaría incompleta la arquitectura de los modelos y la justificación de las estimaciones y de la identificación de factores probablemente asociados.
    \item Ocurrencia ($O$): 6, porque aún no se ha definido una fuente de datos que reúna los biomarcadores, las señales fisiológicas y los eventos de referencia necesarios.
\end{itemize}

\noindent R3 $\rightarrow$ Subestimación del esfuerzo y demoras en la elaboración de los entregables.
\begin{itemize}
    \item Severidad ($S$): 7, porque podría producir una entrega incompleta o reducir el tiempo disponible para revisar el diseño y el análisis de negocio.
    \item Ocurrencia ($O$): 5, porque el trabajo combina química, electrónica, IA y gestión, con dependencias entre las decisiones de cada área y tareas de integración aún por organizar.
\end{itemize}

\noindent R4 $\rightarrow$ Cambios tardíos o ampliación del alcance durante el desarrollo del trabajo.
\begin{itemize}
    \item Severidad ($S$): 6, porque obligaría a rehacer requerimientos, diagramas, análisis de riesgos y estimaciones, afectando los plazos y la coherencia del trabajo.
    \item Ocurrencia ($O$): 4, porque podrían surgir propuestas de agregar biomarcadores o funciones mientras se revisa la bibliografía, aunque el panel y los dos modos de predicción ya están delimitados.
\end{itemize}

\noindent R5 $\rightarrow$ Inconsistencias entre el diseño técnico y las conclusiones de factibilidad y negocio.
\begin{itemize}
    \item Severidad ($S$): 6, porque podría invalidar la trazabilidad de los entregables y llevar a conclusiones apoyadas en prestaciones, costos o supuestos no justificados.
    \item Ocurrencia ($O$): 5, porque las sucesivas revisiones del parche y del alcance pueden dejar versiones diferentes en el objetivo, las tablas, el FODA, el AMFE y el análisis de mercado.
\end{itemize}

    \paragraph{Matriz de Riesgos y Acciones de Mitigación}
    {\consigna
    Aquí debe incorporarse la Matriz de Riesgos y las Acciones Mitigadoras a tomar.\\
    Esta parte se irá completando a lo largo del diseño preliminar y las demás secciones.\\
    El alumno debe investigar diferentes modelos de matrices de riesgos y elegir la que mejor aplique a su proyecto.
    }

\noindent\textbf{c) Tabla de gestión de riesgos}

Se utiliza el índice de prioridad definido para este trabajo como $\mathrm{RPN} = S \times O$, con un umbral máximo aceptable propuesto de 25.

{\small
\renewcommand{\arraystretch}{1.2}
\noindent\begin{tabularx}{\textwidth}{|*{4}{>{\centering\arraybackslash}X|}}
\hline
\textbf{Riesgo} & $S$ & $O$ & \textbf{RPN} \\
\hline
R1 & 9 & 6 & 54 \\
\hline
R2 & 8 & 6 & 48 \\
\hline
R3 & 7 & 5 & 35 \\
\hline
R4 & 6 & 4 & 24 \\
\hline
R5 & 6 & 5 & 30 \\
\hline
\end{tabularx}
\par}

Se proponen acciones específicas para R1, R2, R3 y R5, cuyos RPN superan 25. R4 se acepta provisionalmente con seguimiento, sin un plan específico adicional: cualquier cambio de alcance deberá analizarse antes de incorporarlo, y su puntuación se revisará si comienza a generar retrabajo o demoras.

\noindent\textbf{d) Plan de mitigación de los riesgos que exceden el umbral}

Las acciones se orientan a reducir la posibilidad de desvíos durante la elaboración del trabajo, sin eliminar la incertidumbre técnica ni ampliar los entregables a una validación experimental. Los valores residuales son objetivos estimados sujetos a revisión, no reducciones ya comprobadas.

\noindent R1 $\rightarrow$ Estrategia de mitigación.
\begin{itemize}
    \item Acción: Revisar primero la evidencia de cada zona de sensado y construir una matriz con principio de medición, rango, condiciones de ensayo y limitaciones. Distinguir la medición en sudor de la estimación de glucosa sanguínea. Revisar con los docentes las brechas que afecten al objetivo antes de cerrar el diseño; si no pueden resolverse, declararlas como limitaciones, sin presentar la factibilidad como demostrada.
    \item Severidad residual ($S^{*}$): 9. Se mantiene el impacto si el riesgo ocurre.
    \item Ocurrencia residual objetivo ($O^{*}$): 2, condicionada a la aplicación y revisión de las acciones propuestas.
    \item $\mathrm{RPN}^{*} = 18$.
\end{itemize}

\noindent R2 $\rightarrow$ Estrategia de mitigación.
\begin{itemize}
    \item Acción: Identificar tempranamente los datos y la literatura disponibles, indicando variables, sincronización, etiquetas y condiciones de acceso. Definir la referencia de evaluación de cada salida y especificar por separado la evaluación del objetivo de 30--60 minutos, la del objetivo de 1--8 horas y la identificación de factores probablemente asociados. Resolver además si se requerirán datos de un CGM y cómo se integrarían. Documentar qué datos habría que obtener en una etapa posterior, sin atribuir a datos parciales una validación del sistema completo.
    \item Severidad residual ($S^{*}$): 8. Se mantiene el impacto si el riesgo ocurre.
    \item Ocurrencia residual objetivo ($O^{*}$): 3, condicionada a la aplicación y revisión de las acciones propuestas.
    \item $\mathrm{RPN}^{*} = 24$.
\end{itemize}

\noindent R3 $\rightarrow$ Estrategia de mitigación.
\begin{itemize}
    \item Acción: Dividir el trabajo por entregables, asignar responsables y fechas internas anteriores a la entrega final, y realizar revisiones periódicas de avance. Reservar tiempo para integrar el documento y priorizar los requerimientos, el diseño preliminar y la justificación de factibilidad antes de ampliar el nivel de detalle.
    \item Severidad residual ($S^{*}$): 7. Se mantiene el impacto si el riesgo ocurre.
    \item Ocurrencia residual objetivo ($O^{*}$): 3, condicionada a la aplicación y revisión de las acciones propuestas.
    \item $\mathrm{RPN}^{*} = 21$.
\end{itemize}

\noindent R5 $\rightarrow$ Estrategia de mitigación.
\begin{itemize}
    \item Acción: Mantener una matriz de trazabilidad entre objetivo, requerimientos, componentes y entregables. Hacer una revisión cruzada antes de cerrar cada versión, comprobando la consistencia del panel, los horizontes de estimación y los avisos. Identificar las fuentes y los supuestos de costos, competidores y encuadre regulatorio; distinguir los aspectos confirmados de los pendientes de evaluación.
    \item Severidad residual ($S^{*}$): 6. Se mantiene el impacto si el riesgo ocurre.
    \item Ocurrencia residual objetivo ($O^{*}$): 2, condicionada a la aplicación y revisión de las acciones propuestas.
    \item $\mathrm{RPN}^{*} = 12$.
\end{itemize}

\noindent\textbf{Tabla final}
\par\medskip
{\small
\renewcommand{\arraystretch}{1.2}
\noindent\begin{tabularx}{\textwidth}{|*{7}{>{\centering\arraybackslash}X|}}
\hline
\textbf{Riesgo} & $S$ & $O$ & \textbf{RPN} & $S^{*}$ & $O^{*}$ & $\mathrm{RPN}^{*}$ \\
\hline
R1 & 9 & 6 & 54 & 9 & 2 & 18 \\
\hline
R2 & 8 & 6 & 48 & 8 & 3 & 24 \\
\hline
R3 & 7 & 5 & 35 & 7 & 3 & 21 \\
\hline
R4 & 6 & 4 & 24 & 6 & 4 & 24 \\
\hline
R5 & 6 & 5 & 30 & 6 & 2 & 12 \\
\hline
\end{tabularx}
\par}

Para R4 se conservan los valores iniciales porque se acepta con seguimiento, sin atribuirle una reducción de riesgo. La aceptación y los objetivos residuales deberán reevaluarse en cada revisión de avance.

\subsection{Análisis de mercado y financiero}
{\consigna
En esta subsección se aborda el mercado destinatario del proyecto, la receptividad a la adopción y los aspectos financieros o de soporte relacionados con la realización y éxito del proyecto.
}

\subsubsection{Análisis de mercado y público objetivo}
{\consigna
En este apartado debe describirse de forma cualitativa quiénes son los usuarios finales del proyecto y quiénes son los adquirentes del proyecto, teniendo en cuenta que pueden diferir.\\
También debe describirse la forma en que ese público satisface actualmente las prestaciones que dará el proyecto y cómo lo hace, cuál es el esfuerzo o costo que asume para ello.\\
También debe dimensionarse de alguna forma el mercado potencial.\\
Luego de leer esta sección, el lector deberá tener un entendimiento holístico del usuario final y adquirente: sus necesidades, expectativas, exigencias y motivaciones relacionadas con el proyecto.
}

Para dimensionar el público potencial en Argentina se toma como referencia la población con diabetes, en lugar de extrapolar directamente la facturación mundial del mercado de monitoreo continuo de glucosa. La Federación Internacional de Diabetes estima que en 2024 había aproximadamente 4,3 millones de personas de entre 20 y 79 años con diabetes en Argentina.\footnote{Federación Internacional de Diabetes, \textit{IDF Diabetes Atlas}, 11.\textsuperscript{a} edición, 2025, capítulo 4: estimaciones para América del Sur y Central correspondientes a 2024.} Esta cifra incluye distintos tipos de diabetes y casos no diagnosticados; no representa la cantidad de usuarios de CGM ni de compradores potenciales confirmados del sistema propuesto.

El público objetivo descrito anteriormente está compuesto por personas con diabetes tipo 2 o prediabetes. Como se detalla a continuación, el usuario final previsto es quien ya usa o considera usar un CGM y busca información adicional de contexto. La cifra anterior permite establecer un orden de magnitud de la población con diabetes, pero no cuantifica ese segmento específico ni incluye a las personas con prediabetes. Su tamaño en Argentina queda pendiente de estimar. Para precisar el mercado accesible, podrían relevarse aspectos como el uso de CGM, el acceso a un teléfono compatible, la cobertura de salud y la disposición a utilizar y pagar por un dispositivo complementario; estos aspectos se proponen como líneas de análisis, no como criterios de segmentación ya definidos ni como resultados de un estudio realizado. En esta etapa no se dispone de una estimación propia de ventas o facturación del mercado argentino.

El usuario final es la persona con diabetes tipo 2 o pre-diabetes que ya usa o considera usar un CGM y busca relacionar sus variaciones con información de contexto. El adquirente puede ser el propio paciente, una obra social/prepaga, o una clínica que ofrezca el servicio como parte de un programa de seguimiento. Hoy, este público intenta interpretar esas relaciones mediante consultas, registros manuales de comidas y actividad o aplicaciones de terceros. El sistema busca complementar esas fuentes mediante la identificación de factores probablemente asociados, sin sustituir la evaluación profesional ni establecer causalidad clínica.

\subsubsection{Análisis de competidores y sustitutos}
{\consigna
En este apartado debe describirse de qué otra forma los usuarios finales pueden satisfacer las prestaciones que el proyecto genera.\\
Esto incluye una descripción de competidores existentes y de soluciones sustitutas.\\
Debe contemplarse que no hacer nada también puede ser una opción válida para el usuario final.
}

Los competidores directos son los CGM ya establecidos (Dexcom G7, Libre 3, Accu-Chek SmartGuide), que ofrecen seguimiento de glucosa y distintas funciones de apoyo. La propuesta busca diferenciarse mediante información contextual sobre factores probablemente asociados, sin presentarla como explicación causal. Los sustitutos incluyen el registro manual de comidas y actividad en aplicaciones de terceros, y el seguimiento médico tradicional mediante análisis de laboratorio periódicos. No adoptar un sistema adicional y continuar con controles puntuales también constituye una alternativa para el usuario.

\subsubsection{Preanálisis de costos y precios}
{\consigna
En esta subsección deben analizarse los costos y precios que existen en relación con la realización del proyecto y su oferta futura.
}
    \paragraph{Costos directos e indirectos}
    {\consigna
    En esta subsección deben analizarse los costos que deberá afrontar el proyecto, tanto durante su realización.\\
    Esta sección se actualiza a lo largo del proyecto, manteniendo la línea base y los ajustes necesarios para poder analizar los desvíos que ocurran.\\
    Deben identificarse los costos directos e indirectos, así como indicar si son variables o fijos.
    }
    \paragraph{Precios y receptividad del mercado}
    {\consigna
    En esta subsección deben analizarse los precios de competidores o sustitutos y los valores que el usuario o adquirente está dispuesto a afrontar para disponer del resultado del proyecto.\\
    Para analizar la receptividad del mercado, deben analizarse soluciones que el usuario utiliza que son parte de la problemática o contexto donde el proyecto se desempeña.
   }


\subsubsection{Análisis financiero}
{\consigna
En esta subsección se debe analizar el proyecto desde el punto de vista financiero, es decir, analizar las necesidades económicas del proyecto teniendo en cuenta el momento en que deben ser atendidas. También contempla una proyección de ingresos que el proyecto generaría después de obtener el resultado, en caso de que el proyecto tenga un fin que genere ingresos económicos y un retorno de la inversión económica requerida.\\
Esta es la única parte del presente documento que aborda una perspectiva posterior al proyecto, ya que el retorno de inversión se genera después de que este finaliza. Recomendamos aclarar esta particularidad en este apartado.
}

    \paragraph{Flujo de fondos esperado}
    {\consigna
    En este apartado se debe realizar una proyección estimativa del flujo de fondos esperado para el proyecto, contemplando las dos etapas típicas: la etapa de inversión que ocurre durante el desarrollo del proyecto, donde se visualizan los costos que deben afrontarse y en qué momento; y luego debe incluirse una etapa de explotación económica del resultado (en caso de que el proyecto sea para obtener beneficios económicos), donde se incorporen los ingresos que el proyecto generaría según las expectativas del autor.\\
    La estimación de la generación de ingresos debe estar justificada con datos e información objetiva. El análisis de precios y receptividad del mercado realizado en la sección previa debería contener la información que permita realizar esta proyección si se hizo adecuadamente.\\
    Normalmente el flujo de fondos debería dar negativo en la primera etapa y positivo en la segunda si el proyecto busca un beneficio económico.\\

    {\bf Supuestos y Escenarios posibles}\\
    Como parte del flujo de fondos previsto, deben explicarse aquellos supuestos que dan soporte a la estimación. Algunos ejemplos de los supuestos pueden ser: rango de costos de recursos asumidos, costos de alquileres, carga impositiva prevista, colaboraciones no económicas que generen ahorros o eviten gastos, costos de logística, costos variables de comercialización, marketing, etc.\\
    Normalmente, de no cumplirse un supuesto, los costos previstos van a sufrir variación, generalmente en un mayor costo. Si es posible, pueden dejarse claros estos impactos.\\
    De la misma forma, deben explicarse los supuestos de los ingresos previstos, por ejemplo, precio de venta, cantidad de ventas mensuales, etc.\\
    Posibles variaciones en valores supuestos de costos e ingresos, tomando casos más optimistas o más pesimistas, generan diferentes Escenarios posibles.\\
    En el extremo, los costos e ingresos pesimistas generarán el peor escenario y ambos optimistas el mejor de los escenarios; lo esperable es que el proyecto se ubique en algún punto intermedio y es importante analizar esta realidad.\\
    Previsiones realistas de costos e ingresos darán escenarios probables y la simulación de diferentes escenarios permitirá comprender las posibles realidades del proyecto. Debemos movernos siempre en escenarios realistas.
    }

    \paragraph{Viabilidad financiera: TIR y VAN}
    {\consigna
    En esta subsección se deben analizar los indicadores TIR y VAN que surgen del flujo de fondos esperado, así como proponer una tasa de descuento adecuada para el financiamiento del proyecto.
    }

    \paragraph{Financiamiento necesario y estrategias de soporte}
    {\consigna
    En esta subsección se debe analizar el financiamiento necesario que surge del flujo de fondos y definir, a partir del mismo, una estrategia y alternativas de soporte del flujo de fondos.\\
    Pueden definirse diferentes alternativas para poder soportar el proyecto en caso de que el financiamiento que se obtenga no sea suficiente, qué recortes pueden realizarse o qué tácticas pueden emplearse para disminuir las necesidades de financiamiento.
    }